\newpage
\section{Conclusions}\label{sec:conclusion}

This paper asked when a large-cap, long-only momentum portfolio is worth tilting away from its benchmark, and answered it with a two-stage framework: a drawdown-only probabilistic regime signal routes a nonlinear cross-sectional model between classical winner selection in calm markets and term-structure-wide reversal at panic peaks. The honest headline is that the strategy earns no unconditional risk-adjusted alpha in this configuration; its value is conditional and implementational. Essentially all benchmark-beating return arrives in the quarter of months the regime signal classifies as panic, and the signal's practical product is a when-to-trade rule that concentrates turnover where the payoff is.

\subsection{Conclusions and Contributions}\label{sec:conclusions_contributions}
Cross-sectional momentum's regime-dependent structure is real and admits a concrete characterization. The regime signal acts as a context variable, conditioning how a nonlinear cross-sectional model uses the full momentum term structure rather than predicting returns directly: it carries roughly five times its equal-weight share of the model's attribution, and once it saturates in panic the momentum features take over the stock-by-stock differentiation (Section~\ref{sec:horizon_evidence}). The analysis makes three contributions.

First, the methodological combination of a continuous probabilistic regime filter with a nonlinear cross-sectional model remains novel relative to prior regime-conditional momentum work, which paired simpler regime treatments (binary indicators, threshold-based blends, or fast-slow weights) with linear or rank-based cross-sectional rules. The applied configuration sharpens the recipe to its parsimonious core: a univariate hidden Markov model on market drawdown, chosen on literature grounds and confirmed by a train-only selection rule that never adds a second stress feature (Section~\ref{sec:feature_selection}).

Second, the empirical analysis characterizes how the cross-section reorganizes across regimes at the level of individual stock-month picks. A four-cluster decomposition of the model's own selections traces the cycle from strong winner selection in calm ($+1.2\,\sigma$ above the cross-sectional mean in the continuation mode) to term-structure-wide reversal at panic peaks ($-1.0\,\sigma$ below, negative at every one of the twelve horizons), a composition no single-horizon rank-based strategy can produce. \citet{Cooper2004} noted this state-dependent asymmetry descriptively from a binary backward-looking market state; this paper implements the regime-conditioned response as a single nonlinear model with a fully-invested long-only book, composition rather than exposure absorbing the regime change.

Third, and most useful to a practitioner, the regime signal converts into an execution rule. The strategy's active return is panic-concentrated while roughly three-quarters of a monthly-rebalanced book's trading occurs in calm months that earn nothing. A fixed trading band alone raises the net information ratio at 10\,bp costs from $0.06$ to $0.29$; conditioning the band on the regime, wide in calm and tight in panic, extends the frontier to $0.39$ and the breakeven cost from $14$ to $42$\,bp, a lead over the best static band that holds at every cost level but whose final increment is not individually significant (Section~\ref{sec:cost}). To our knowledge no prior work implements a state-dependent buy/hold band on a cross-sectional selection book: the theory literature anticipates state-dependent no-trade regions \citep{Jang2007, LynchTan2011} but has not taken them to portfolio data, and the closest empirical antecedent conditions trading speed, not a no-trade band, on the regime \citep{CollinDufresne2020}.
The companion finding is that suppressing rank-driven selling in panic is, in point estimate, not merely cheaper but better gross of costs (information ratio $0.21$ to $0.33$). Together with the calm-month evidence that the signal keeps the book benchmark-like when the tilt earns nothing, the regime signal's deepest contribution is discipline about when not to trade.

\subsection{Limitations}\label{sec:limitations}

\paragraph{Reversal-failure vulnerability and test-window scope}
The strategy's primary vulnerability is a prolonged economic deterioration in which the prior decline reflects genuine business failure rather than temporary mispricing, so the beaten-down stocks the model selects do not recover. The panic-mode bet works in short panics but fails in a sustained bear: in the 2000--2002 dot-com bust the extended backtest loses $65.6\%$ against the benchmark's $36.4\%$, setting the full-sample maximum drawdown of $77.3\%$ (Section~\ref{sec:historical_stress}). The regime classification was correct in that episode; what changed was the panic-regime pattern, with businesses failing rather than prices being temporarily depressed, and the model cannot distinguish these cases ex ante. The 2011--2025 production window avoids this failure mode, and the 2000--2025 extension documents it exactly once, so we cannot quantify how the strategy would behave across the distribution of such episodes.

A test window containing a further prolonged bear market would deliver a different point estimate. The exposure-scaling overlay discussed in Section~\ref{sec:future_work} addresses the structural vulnerability as a circuit breaker. A depression-style scenario lies outside the 25-year backtest, and it is hard to assess via simulation: an accurate stress test would have to model the heterogeneous business-failure dynamics that drive the failure mode, not just shock all stock returns uniformly by a chosen magnitude over a sustained period.

\bigskip\paragraph{Regime stationarity under an expanding window}
The regime model is re-estimated each year, but on an expanding window in which the 1990s and 2000s permanently dominate the training sample. If the market undergoes structural change (a persistent shift in baseline volatility from algorithmic trading, or a different monetary-policy regime), thresholds anchored to decades-old drawdown behavior may miscalibrate, classifying structurally calm periods as panic. The sustained elevation of $\pi_t^{\text{filter}}$ in parts of the recent sample may partly reflect such drift rather than genuine prolonged stress. A rolling estimation window, or online updating that discounts distant history, would address this directly, at the cost of noisier parameters.

\bigskip\paragraph{In-sample choices and statistical power}
The regime indicator itself is protected against specification search: drawdown is fixed on literature grounds and the data-driven check uses training information only (Section~\ref{sec:feature_selection}). The residual in-sample choices sit elsewhere. The execution-band grid is selected on the full 179-month window, the candidate pool behind the selection rule was assembled with knowledge of pre-2011 research, and the net value of every execution policy is concentrated in the 2018--2025 half of the sample. Statistical power is the binding constraint throughout: with 45 panic months, the panic-subsample alphas are economically large under some classifications and insignificant under all of them, and the information-ratio differences across ablation variants are directional rather than individually significant. The claims the paper rests on are the ones that clear inference in their exact form: the banded implementation beats monthly rebalancing of the same signal ($t = 2.99$), and the tighter-in-panic direction beats its mirror in all ten pairs (sign test $p = 0.001$). The regime band's lead over the best static band, the panic-subsample alphas, and the claim of beating the index outright ($t = 1.40$) do not clear inference; the panic concentration of active return is economically unambiguous but is documented without a formal subsample test.

\subsection{Recommendations for Future Research}\label{sec:future_work}

Five directions extend this framework most directly.

\paragraph{Richer regime-signal representation}
The current model uses only the contemporaneous filtered probability $\pi_t^{\text{filter}}$, discarding regime dynamics. Two complementary extensions follow. The first looks forward, computing a $k$-step regime forecast $P(s_{t+k} = 1 \mid \mathbf{z}_{1:t})$ from the HMM's transition probabilities so the portfolio can reposition before a regime shift rather than after it. The gain is concentrated in transition months, where the current framework trades on a stale view of the regime. The second looks backward, feeding the trajectory of $\pi_t^{\text{filter}}$ rather than its current value so that a fresh spike (where mean reversion is most likely) and a sustained elevation (where the reversal mechanism may be breaking down) become distinguishable. A sequential model, or simple lag-based feature engineering (lagged $\pi_t^{\text{filter}}$ values, a short-window rate of change, a regime-duration counter), would expose this trajectory. Together, the two extensions provide anticipation and memory, and the trajectory features bear directly on the reversal-failure limitation above.

\bigskip\paragraph{Alternative nonlinear models}
The selection rule found here is what XGBoost can learn from the regime probability and the twelve momentum horizons; a more expressive model may recover a sharper version of the same signal. The most natural alternative is a feedforward neural network, which \citet{Gu2020} find competitive with tree ensembles for cross-sectional return prediction. Beyond raw performance, the robustness question is the interesting one: if the regime-conditional selection shift recurs in a non-tree architecture, the mechanism is a property of the data rather than an artefact of tree splits, and comparing the network's per-feature attribution to the SHAP decomposition in Table~\ref{tab:shap} would reveal whether the regime signal plays the same routing role outside the tree representation.

\bigskip\paragraph{Beyond the long-only mandate}
The applied study deliberately excludes the short side, where large-cap borrow is cheap but the crash risk of shorting panic-month rebounds is the classic momentum failure mode. The short side need not be a mirror image of the long leg. Whether a modest short overlay adds value net of borrow costs and crash risk in a large-cap universe, and whether the regime signal should gate the short book more aggressively than the long book, are open implementation questions.

\bigskip\paragraph{Regime-conditional feature sets}
Section~\ref{sec:fund_ablation} shows that adding firm-level fundamentals inside the regime-conditioned model destroys its panic-month active return while a momentum-plus-fundamentals model without the regime signal remains competitive, suggesting the optimal feature set is itself regime-dependent. A layered construction in which the model produces the ranking and fundamentals enter as a post-ranking quality filter on the long book is a natural starting point; a more ambitious variant trains separate cross-sectional models for the calm and panic branches, each with its own feature set. Both designs require careful out-of-sample validation because the regime-conditional split reduces effective sample size.

\bigskip\paragraph{Exposure-scaling overlay}
An exposure-scaling overlay in the spirit of \citet{DanielMoskowitz2016}, layered on top of the stock selection, would provide a circuit breaker against sustained downturns that violate the panic-rebound premise. The overlay requires a triggering rule; candidates include $\pi_t^{\text{filter}}$ above a threshold for a minimum number of consecutive months, or its rolling maximum over a fixed window exceeding a calibrated bound. The 2011--2025 production window cannot discriminate among candidate rules, because it contains no prolonged bear market and any overlay only reduces exposure in months that turn out to be profitable; the dot-com sub-period of the 2000--2025 extension (Section~\ref{sec:historical_stress}) is the appropriate setting to evaluate them, and a rule that converts even part of that episode's $-65.6\%$ into benchmark-like performance would remove the strategy's worst-case objection.
